\documentclass[letterpaper,10pt]{article}

% One-page base (from Ryan's one-page PDF). Same layout as resume.tex.
% Compiles with pdfLaTeX, including on Overleaf's default settings.
\usepackage{iftex}
\ifPDFTeX
  \usepackage[T1]{fontenc}
  \usepackage[utf8]{inputenc}
  \usepackage[scaled=0.92]{helvet}
\else
  \usepackage{fontspec}
  \setsansfont{texgyreheros}[Extension=.otf, UprightFont=*-regular, BoldFont=*-bold, ItalicFont=*-italic, BoldItalicFont=*-bolditalic, Scale=0.92]
\fi
\renewcommand{\familydefault}{\sfdefault}
\usepackage[letterpaper,margin=0.5in]{geometry}
\usepackage{enumitem}
\usepackage{titlesec}
\usepackage[hidelinks]{hyperref}

% Makes the PDF text map cleanly to Unicode for ATS parsers.
\ifdefined\pdfgentounicode
  \input{glyphtounicode}
  \pdfgentounicode=1
\fi

\pagestyle{empty}
% No line-end hyphenation, so keywords stay whole for ATS parsers.
\hyphenpenalty=10000
\exhyphenpenalty=10000
\emergencystretch=2em
\setlength{\parindent}{0pt}
\setlength{\parskip}{0pt}
\setlist[itemize]{nosep,leftmargin=1.3em,itemsep=0.5pt,topsep=1pt}
\titleformat{\section}{\large\bfseries}{}{0em}{}[\titlerule]
\titlespacing{\section}{0pt}{6pt}{3pt}

\newcommand{\role}[3]{\textbf{#1} \hfill #3\par\textit{#2}\par}

\begin{document}

\begin{center}
{\LARGE\textbf{Ryan McPhedrain}}\par\vspace{4pt}
{\normalsize Neuroscientist and Postdoctoral Researcher}\par\vspace{3pt}
\small
\href{mailto:rmcphedr@gmail.com}{rmcphedr@gmail.com} | 236-512-2798 | \href{https://linkedin.com/in/ryanmcphedrain}{linkedin.com/in/ryanmcphedrain} | \href{https://github.com/rmcphedr}{github.com/rmcphedr}
\end{center}

\section{Summary}
Neuroscientist and Mila postdoctoral researcher with 9+ years designing hypothesis-driven studies and analytical pipelines for complex biological data. Combines electrophysiology, two-photon imaging, and molecular neuroscience with computational neuroscience, machine learning, and multimodal integration. Experienced coordinating rigorous research, mentoring trainees, and translating scientific questions into reproducible analyses and publications.

\section{Skills}
\textbf{Experimental Neuroscience:} Patch-Clamp Electrophysiology, Two-Photon and Fluorescence Imaging, Synaptic Development and Neural Circuits, Molecular Biology, Cell Culture\par\vspace{2pt}
\textbf{Computational Neuroscience and AI:} Python (NumPy, Pandas, scikit-learn, PyTorch, MNE, Nilearn), Statistical Modeling, Machine Learning, Multivariate Analysis (CCA, PLS), Multimodal Integration\par\vspace{2pt}
\textbf{Scientific Leadership:} Experimental Design, Research Strategy, Scientific Writing and Grant Preparation, Cross-Disciplinary Collaboration, Mentorship

\section{Experience}
\role{Postdoctoral Researcher}{Mila, Montreal, QC}{2024 -- Present}
\begin{itemize}
    \item Designed reproducible Python pipelines to evaluate competing fMRI denoising and statistical strategies across large, multi-study datasets, supporting quality control, hypothesis testing, and figures for research published in Nature Medicine.
    \item Coordinated neuroscientists, machine-learning researchers, statisticians, and international collaborators across 3 continents to harmonize 13 studies, 300+ participants, and 500+ scans, enabling a cross-study analysis of psychedelic effects on brain functional connectivity.
    \item Developed multimodal analyses integrating rodent transcriptomics with longitudinal human neuroimaging to investigate conserved relationships between spatial gene expression and brain activity.
    \item Applied machine learning and BERT-based semantic analysis to complex behavioral data; contributed to manuscripts and grant proposals and mentored graduate trainees in statistics, programming, and scientific communication.
\end{itemize}
\vspace{3pt}

\role{Research Consultant (Internship)}{The Decision Lab, Montreal, QC}{2023 (4 months)}
\begin{itemize}
    \item Built an interactive NLP application to analyze demographic differences in financial attitudes, translating complex behavioral data into clear outputs for multidisciplinary stakeholders.
    \item Used PCA and decision trees to identify behavioral subgroups in a smoking-cessation program and inform personalized interventions under project constraints.
\end{itemize}
\vspace{3pt}

\role{Graduate Researcher}{McGill University, Montreal, QC}{2016 -- 2023}
\begin{itemize}
    \item Designed and executed hypothesis-driven research on synaptic development and neural-circuit formation using in vivo two-photon imaging, ex vivo patch-clamp electrophysiology, CRISPR-based methods, molecular biology, and statistical analysis.
    \item Developed reproducible Python image-analysis tooling, authored peer-reviewed research, presented at scientific conferences, and trained graduate researchers in experimental methods and analysis.
\end{itemize}

\section{Education}
\role{PhD in Neuroscience (4.0 GPA)}{McGill University, Montreal, QC}{2016 -- 2023}
\role{BSc (Honours) in Molecular Biology \& Biochemistry (3.93 GPA)}{Simon Fraser University, Burnaby, BC}{2010 -- 2015}

\end{document}
